% Camino de senal del amplificador de ECG, dibujado para una columna IEEE.
\begin{circuitikz}[font=\footnotesize]
\ctikzset{bipoles/length=0.85cm, amplifiers/scale=0.62, diodes/scale=0.5,
  resistors/scale=0.7, capacitors/scale=0.6, grounds/scale=0.6}
\tikzset{dot/.style={circle, fill, inner sep=0pt, minimum size=2.6pt}}
% ---------------- Etapa 1: amplificador de instrumentacion ----------------
\node[op amp, noinv input up] (B) at (3.0,3.0) {};
\node[op amp] (A) at (3.0,0.2) {};
\node at (3.25,3.6) {U1B};
\node at (3.25,-0.4) {U1A};
% entrada + (LA): R1 y diodos D1/D3 a masa
\draw (0,0 |- B.+) node[ocirc] (la) {} to[R, l=$R_1$, a=10k] (0.85,0 |- B.+)
  -- (B.+);
\node[left] at (la) {LA};
\draw (1.0,0 |- B.+) node[dot] {} to[D] ++(0,-0.55) node[ground] {};
\draw (1.27,0 |- B.+) node[dot] {} to[D, invert] ++(0,-0.55) node[ground] {};
% entrada - (RA): R2 y diodos D2/D4 a masa
\draw (0,0 |- A.+) node[ocirc] (ra) {} to[R, l_=$R_2$, a^=10k] (0.85,0 |- A.+)
  -- (A.+);
\node[left] at (ra) {RA};
\draw (1.0,0 |- A.+) node[dot] {} to[D] ++(0,-0.55) node[ground] {};
\draw (1.27,0 |- A.+) node[dot] {} to[D, invert] ++(0,-0.55) node[ground] {};
% R3 entre las entradas inversoras
\draw (B.-) -- (1.7,0 |- B.-) node[dot] {} to[R, l_={$R_3$ 10k}] (1.7,0 |- A.-) node[dot] {} -- (A.-);
% realimentaciones R5 y R4
\draw (1.95,0 |- B.-) node[dot] {} |- (2.25,2.05) to[R, l={$R_5$ 499k}] (3.4,2.05) -| (3.95,0 |- B.out) -- (B.out);
\draw (1.95,0 |- A.-) node[dot] {} |- (2.25,1.15) to[R, l_={$R_4$ 499k}] (3.4,1.15) -| (3.95,0 |- A.out) -- (A.out);
% ---------------- Etapa 2: restador con filtro pasabanda ----------------
\node[op amp] (C) at (6.95,1.6) {};
\node at (7.25,1.0) {U1C};
\draw (3.95,0 |- B.out) node[dot] {} to[R, l=$R_6$, a=100k] (4.85,0 |- B.out)
  to[C, l=$C_3$, a=33$\mu$] (5.6,0 |- B.out) |- (C.-);
\draw (3.95,0 |- A.out) node[dot] {} to[R, l_=$R_7$, a^=100k] (4.85,0 |- A.out)
  to[C, l_=$C_4$, a^=33$\mu$] (5.6,0 |- A.out) |- (C.+);
% realimentacion R10 || C7
\draw (5.85,0 |- C.-) node[dot] {} -- (5.85,3.3) node[dot] {} to[R, l_={$R_{10}$ 100k}] (7.95,3.3) node[dot] {}
  -- (7.95,0 |- C.out) node[dot] {} -- (C.out);
\draw (5.85,3.3) -- (5.85,4.0) to[C, l={$C_7$ 15n}] (7.95,4.0) -- (7.95,3.3);
% rama a masa de la entrada +: C8 y R8 + R9 + RV1
\draw (5.85,0 |- C.+) node[dot] {} -- (5.85,0.05) node[dot] {} to[C, l_={$C_8$ 15n}] (5.85,-0.75) node[ground] {};
\draw (5.85,0.05) -- (6.95,0.05) to[R, l={$R_8{+}R_9$ 98k}] (6.95,-0.85) to[vR, l={RV1 5k}] (6.95,-1.75) node[ground] {};
% ---------------- Etapa 3: inversor de salida ----------------
\node[op amp] (D) at (10.75,1.29) {};
\node at (11.05,0.68) {U1D};
\draw (7.95,0 |- C.out) to[R, l=$R_{11}$, a={4,7k}] (8.95,0 |- C.out)
  to[vR, l=RV2, a=10k] (9.8,0 |- C.out) node[dot] {} -- (D.-);
\draw (D.+) -- (9.8,0 |- D.+) node[ground] {};
\draw (9.8,0 |- C.out) -- (9.8,2.45) to[R, l={$R_{12}$ 100k}] (11.7,2.45) -- (11.7,0 |- D.out) node[dot] {} -- (D.out);
\draw (11.7,0 |- D.out) -- (12.0,0 |- D.out) node[ocirc] (vo) {};
\node[right] at (vo) {$v_o$};
\end{circuitikz}
